\documentclass[12pt,a4paper]{article}
\usepackage{fontspec}
\IfFontExistsTF{Times New Roman}{\setmainfont{Times New Roman}}{\setmainfont{DejaVu Serif}}
\usepackage[russian]{babel}
\usepackage{amsmath,amssymb,booktabs}
\usepackage[a4paper,margin=23mm,top=20mm,bottom=22mm]{geometry}
\usepackage{xcolor}
\definecolor{ink}{HTML}{193A4D}
\usepackage[unicode,colorlinks=true,urlcolor=ink,linkcolor=ink]{hyperref}
\setlength{\emergencystretch}{2em}
\setlength{\parindent}{1.1em}
\setlength{\parskip}{0.14em}
\raggedbottom
\hypersetup{pdftitle={Рецензия на доказательство о числе положительных решений}}\begin{document}
\begin{center}{\Large\bfseries\color{ink} Рецензия на доказательство о числе положительных решений\par}\end{center}

24 ноября 2026 года. Предмет проверки — рукопись GPT Pro, предоставленная пользователем, §§1–9. В рецензии сопоставлены аналитический аргумент, точные вычисления и проверенные ядром Lean утверждения.

\section*{\color{ink} Заключение}

Для цепи длины \(n\ge4\) с входом в камеру 1, измерением количеств в камерах \(3,n\) и одной известной точкой выведения \(\ell\) представленное доказательство даёт точное правило для всех длин. Наибольшее число различных положительных \textbf{регулярных} наборов скоростей, порождающих одни и те же данные, равно одному при \(\ell=3\) или \(\ell=n\) и двум при любом другом \(\ell\). Обе границы достигаются. Ошибки, которая меняла бы этот ответ, при независимой проверке не обнаружено.

Регулярность означает, что определитель якобиана квадратной карты наблюдений ненулевой. Два регулярных прообраза сохраняются на общем открытом множестве данных. Число два в теореме — максимум. Для конкретных данных число положительных решений может зависеть от данных, а вырожденные точки в определение не входят. Теорема относится к точным измерениям и фиксированному положению входа и выходов.

\section*{\color{ink} Почему верхняя граница работает}

Последний выход имеет передаточную функцию \(q/D(s)\) с \(q>0\), поэтому однозначно фиксирует монический знаменатель и \(q\). Другой выход фиксирует \(pW(s)\), а его старший коэффициент равен \(p>0\). Это даёт эквивалентность исходного эксперимента и наблюдения \(F=(D,W,p,q)\) даже при возможном сокращении первой дроби.

При взаимно простых \(D,W\) спектральные веса последней камеры вычисляются по квадратному многочлену \(P\) из формулы (1). Нормировка весов даёт аффинную плоскость. Условие на \(p\) при \(\ell\ge3\) либо условие на \(K=qW(0)/p\) при \(\ell\le2\) оставляет прямую в этой плоскости у регулярного решения. Если последнее условие зависит от нормировки и оставляет всю плоскость, то на ней ранг двух оставшихся наблюдений не превосходит одного.

Для всех положительных точек на прямой выполняется \(q/D(0)=\Phi_{n-\ell}\). Веса входят линейно в энергию каждого допустимого многочлена. Минимум этих энергий вогнут по весам. Если три упорядоченные точки имеют одно и то же значение, левый и правый секущие наклоны в средней точке равны нулю. Производная в ней также равна нулю. Это противоречит регулярности. При \(\ell=3\) и \(\ell=n\) оба оставшихся ограничения линейны в коэффициентах \(P\). Две различные точки тогда породили бы отрезок решений, несовместимый с локальной обратимостью в его концах.

Данные с общим корнем \(D,W\) тоже охвачены. Конечное число различных регулярных прообразов продолжается на общую открытую окрестность данных по теореме об обратной функции. В ней есть точка с ненулевым результантом. Любое нарушение границы перенеслось бы в уже разобранный случай.

\section*{\color{ink} Почему граница достигается при любой длине}

При \(\ell=1,2\) семейство §5 даёт положительные веса внутри интервала. В каждом из двух пределов остаются \(n-2\) узлов. Многочлен степени \(n-2\) с единицей в нуле зануляется на каждом таком носителе, поэтому на обоих концах энергия стремится к нулю. Внутри интервала минимум достигается и положителен. Ненулевой поперечный дифференциал \(K_v\) даёт полный ранг в двух пересечениях подходящей горизонтали. Для \(\ell=3,n\) формулы производных §6 дают по одному регулярному примеру.

Для \(4\le\ell<n\) база \((5,4)\) имеет значения \(\Phi\), равные \(1/4,\ 4/15,\ 1/4\), в трёх точках одного интервала. Оба перехода §8 сохраняют строгое превышение среднего значения над крайними. В переходе A вклад нового узла в старший нужный момент имеет порядок \(O(L^{5-2\ell})\). В переходе B сначала выбирается одно малое \(\varepsilon>0\), затем одно достаточно большое конечное \(L\) для всех трёх точек. В пределе множитель \(S(t)/(S(t)+\varepsilon)\) зависит от точки. Именно выбор \(\varepsilon\) до выбора \(L\) сохраняет оба строгих неравенства. Положительность на всём интервале следует из аффинности ненормированных весов по \(t\). Из базы достаточно выполнить \(\ell-4\) переходов A и \(n-\ell-1\) переходов B. Это доказательство достижимости для произвольных длин.

\section*{\color{ink} Уточнения, внесённые в текст}

1. \textbf{§8, строки 835–853 исходного текста.} Для нижней предельной оценки сначала выбирается последовательность \(L\to\infty\), реализующая \(\liminf\) энергий. Из неё извлекается подпоследовательность минимизаторов со сходящимися коэффициентами. Ограничение \(r_L(L)=O(\sqrt L)\) обнуляет в пределе старший коэффициент. Это доказывает требуемую оценку полного \(\liminf\).

2. \textbf{§9, строка 899.} Определитель замены новых координат на прежние равен \(-12p\). Из перехода от коэффициентов \(D\) к значениям \(D(0),\ldots,D(3)\) получается \(12\). Из перехода \((W(0),p)\mapsto(p,pW(0))\) получается \(-p\). Для старой контрольной точки \((1,2,3,4,5,6,7)\) воспроизведены все четыре определителя \(-311040,-198720,-103680,-51840\).

3. \textbf{§§2, 4, 7, строки 205, 329 и 623–625.} Уточнены знак при смене спектральной переменной, независимость однородного фиксирующего условия от нормировки в совместной системе и ненулевая производная инъективной дробно-линейной параметризации. Последнее обосновывает перенос одномерной регулярности на оставшуюся пару наблюдений в локальных спектральных координатах.

4. \textbf{Лемма о спектральной реконструкции, §2.} Для удлинения цепи важно уметь установить достижимость вновь построенных \(D,W\). Поэтому лемма в английской статье сформулирована для любых взаимно простых монических \(D,W\) нужных степеней с простыми отрицательными корнями \(D\). Положительное нормированное решение уравнений для весов даёт физическую цепь. Эта достаточность уже следует из построения ортогональных многочленов. Её явная формулировка устраняет возможность круговой ссылки при применении к новому спектру.

\section*{\color{ink} Что принял Lean}

Обозначение \(M(n,\ell)\) относится к той же карте положительных физических скоростей, что и в рукописи. В коде \(m=n-1\), а камеры нумеруются с нуля. Поэтому, например, индекс 3 в пятикамерном сертификате означает четвёртую камеру.

\par\noindent\(\bullet\)\enspace Для всех длин установлена эквивалентность равенства физической карты коэффициентов и совпадения обеих выходных кривых для всех непрерывных входов. Проверены три кофакторных тождества §1 и невырожденность базиса Крылова на соответствующем крайнем блоке.
\par\noindent\(\bullet\)\enspace Для всех допустимых размеров проверены вогнутость \(\Phi\), формула для минимума через обратную матрицу Грама, граница в два регулярных пересечения в спектральной прямой и продолжение конечного числа регулярных прообразов в плотное множество данных. Для произвольной положительно определённой матрицы Якоби доказано равенство \(\Phi_m\) разности крайних элементов обратных матриц полной цепи и главного блока. В физической форме \(q/D(0)=\Phi_{n-\ell}\) проверено при предъявлении обращённой матрицы Якоби и совпадении её главных миноров с \(D,H,R,C\). Проверены пределы A и B для сформулированных в коде весов и преобразованных базисов, а также индукция по положениям при наличии двух переходов.
\par\noindent\(\bullet\)\enspace Полностью доказано \(M(4,2)=2\) для физических скоростей. Верхняя оценка в этом частном случае исключает даже три различных положительных решения без предположения их регулярности. Для пяти камер доказано \(M(5,4)\ge2\). У двух различных положительных векторов из §7 совпадают кривые, а определители вещественных производных Фреше равны \(3456\) и \(-4698\).
\par\noindent\(\bullet\)\enspace Общая формула аккуратно записана в Lean как высказывание SharpCount. Для неё в комплекте отсутствует сквозное доказательство. Главная недостающая часть — гладкая обратимая замена всех физических скоростей на спектральные координаты \((D,W,e,f)\), с доказательством, что каждый допустимый вес реализует именно заданный хвост \(W\). Далее требуется связать регулярность полной карты с поперечными производными, реализовать семейства §§5–6 и применить пределы §8 к физическим вычетам во всех длинах. Предел для абстрактных весов или проверенный начальный пример сам по себе этих связей не устанавливает.

\section*{\color{ink} Воспроизведение и источники}

В комплекте находятся исходники Lean 4.19.0 с закреплённой версией mathlib, независимая проверка на стандартном Python и журналы запуска. Вычисления используют рациональные числа. Проверены вариационная формула на 30 физических моделях, симметричное семейство для \(4\le n\le12\), обе операции удлинения и их композиции до \((8,5)\). В пятикамерном примере совпал весь наблюдаемый вектор \((24,150,206,96,17,8,7,2,6)\). Численные проверки служат независимым контролем знаков, индексов и якобианов.

\href{https://arxiv.org/html/2602.20894v1#S3.Thmtheorem1}{Castillo и Gordillo-Núñez, леммы 2.1, 3.1 и теорема 3.1} дают реконструкцию по двум спектрам. \href{https://arxiv.org/html/2507.04496v1#S5}{Meshkat и Shiu, вопрос 5.3} обсуждают степени идентифицируемости камерных моделей. Ни один из этих двух результатов не вычисляет здесь максимум регулярных положительных прообразов с дополнительными ограничениями \(p,q\). Приоритет по отношению ко всей литературе этой рецензией не установлен. Рецензию и формальные проверки выполнили агенты Codex. Внешняя экспертная рецензия остаётся следующим этапом.\end{document}
